\section{Aggregate Lower Bound Theorem}
\label{sec:lower_bound}

\subsection{Realized capability volume ($\volR$)}

Before stating the aggregate theorem, we make the object being
bounded precise.  $\volR$ is the poset measure evaluated on the
\emph{exercised sub-poset}: the restriction of the capability
space to capabilities that are currently being exercised.

\begin{definition}[Exercised Sub-Poset and Realized Capability
Volume]
\label{def:volR}
Let $P$ be the full capability poset with measure
$\volL$~\citep[Proposition~1]{lasser2026poset}.

\medskip\noindent
\textbf{Instantaneous form.}
For each time $t$, the \emph{exercised sub-poset}
$P^{\mathrm{ex}}_t \subseteq P$ is the set of capabilities being
exercised at time $t$---a \emph{latent structural fact} about
the system, independent of any observation mechanism:
\[
  P^{\mathrm{ex}}_t \;=\; \{d \in P : d \text{ is exercised at
  time } t\},
\]
where ``exercised at $t$'' is formalized via the counterfactual
exercise condition of Appendix~\ref{app:exercise_indicator}
(Definition~\ref{def:exercise_indicator}): $d$ is exercised at
$t$ iff removing $d$ at that time would make some realized
cooperative output unrealizable.  The \emph{instantaneous
realized capability volume} is the poset measure restricted to
this sub-poset:
\[
  \volR_t(S) \;=\; \volL\bigl(S \cap P^{\mathrm{ex}}_t\bigr)
  \qquad \text{for any } S \subseteq P.
\]

\medskip\noindent
\textbf{Window-active form (union over time, not a time
integral).}
For a trade window $W = [t_0, t_0 + T]$
(Definition~\ref{def:trade_window}), the \emph{window-active
exercised sub-poset} is the union of the instantaneous
exercised sub-posets across the window,
\[
  P^{\mathrm{ex}}_W \;=\; \bigcup_{t \in W}
  P^{\mathrm{ex}}_t,
\]
and the corresponding \emph{window-active realized
capability volume} is
\[
  \volR^{[W]}(S) \;=\; \volL\bigl(S \cap P^{\mathrm{ex}}_W\bigr)
  \qquad \text{for any } S \subseteq P.
\]
$\volR^{[W]}$ counts a capability as contributing to realized
volume if it was exercised at \emph{any} time during $W$; the
instantaneous form $\volR_t$ counts only capabilities exercised
exactly at $t$.  Both inherit the measure structure of $\volL$
restricted to the corresponding sub-poset, and the
axiom-analogue analysis of
Section~\ref{sec:axiom_inheritance} applies to each.

\medskip\noindent
\textbf{Terminology: ``window-active,'' not a time integral.}
We use the term \emph{window-active} for $\volR^{[W]}$ because
the operation is the union over time
$P^{\mathrm{ex}}_W = \bigcup_{t \in W} P^{\mathrm{ex}}_t$, not
a duration- or frequency-weighted time integral.  A capability
exercised once for an instant during $W$ contributes the same
$\volL$ weight to $\volR^{[W]}$ as a capability exercised
continuously throughout $W$; the diagnostic semantics is
``did the capability witness any exercise in this window?''
not ``how much exercise per unit time accumulated?''
The bracket notation $[W]$ denotes ``indexed by $W$,'' not
``integrated over $W$''; we keep the bracket for symbol
continuity with the broader poset-measurement
literature~\citep{lasser2026scm} but spell out the witness
semantics here to avoid the dimensional misreading that
``integrated'' would invite.  Where duration- or
frequency-weighted exercise is needed (e.g., to distinguish
intermittent from persistent exercise), an explicit weighted
integral is required, which is outside this paper's scope.

\medskip\noindent
\textbf{Convention.}
Unless otherwise noted, $\volR$ without a time subscript refers
to the window-active form $\volR^{[W]}$ for the ambient
trade window $W$.  This is the quantity the sacrifice channel
(Theorem~\ref{thm:aggregate_bound}) actually bounds: a sacrifice
event at $t_n$ records an acquisition, not an instantaneous
exercise witness, so the theorem's lower bound is stated
against $\volR^{[W]}$---which captures any exercise of $Y_n$
anywhere in $W$---rather than against the instantaneous
$\volR_{t}$ at a particular evaluation time.  Assumption~S0
(Section~\ref{sec:assumptions}) is correspondingly stated
directly in window-active form:
$\Uin{n}(Y_n) \leq \volR^{[W]}(Y_n)$.
\end{definition}

\begin{remark}[$\volR$ is latent; $\volRlower$ is the bound]
\label{rem:volR_diagnostic}
$\volR$ (in either form) is defined structurally as a
ground-truth quantity about the system; computing it
\emph{directly} requires evaluating the counterfactual exercise
condition for every capability at every relevant time, which the
Introduction rejects as intractable and privacy-incompatible at
scale.  The revealed-sacrifice observation channel does not
observe $P^{\mathrm{ex}}_t$ or $P^{\mathrm{ex}}_W$ directly; it
produces a lower bound $\volRlower$ on the window-active
$\volR^{[W]}$ restricted to an observed sub-space
$U \subseteq P$.  Inside the observation perimeter $O$ of
Appendix~\ref{app:exercise_indicator}, direct evaluation of
$e_t(d)$ is permitted (no privacy barrier within consented
operational scope) and contributes to a combined diagnostic
$\volRcomb = \volRexact(O) + \volRlower(P \setminus O)$, which
is a valid lower bound on $\volR(P)$ when $O$ is
boundary-closed from $P \setminus O$ in the full
poset-independence sense (neither cooperative-composition nor
subsumption links cross the boundary; see
Definition~\ref{def:exercise_indicator} conditions BC1/BC2).
Outside that regime, the compositional lower-bound
interpretation of $\volRcomb$ is not established by this paper
and the quantity should be treated as a heuristic composite
(Appendix~\ref{app:exercise_indicator},
Remark~\ref{rem:boundary_cooperative_undercount}).  $\volR$ itself is the
latent target; $\volRlower$ (and $\volRexact$ on $O$) are the
observation-channel outputs the framework actually computes.
\end{remark}

\subsection{Event-local decomposition}

\begin{definition}[Event-Local Bundle Decomposition]
\label{def:bundle_decomp}
For each revealed-sacrifice event $(i_n, X_n, Y_n, t_n)$ whose
bundle $Y_n$ is \emph{within-bundle poset-independent} (no
subsumption relations and no cooperative-composition links among
the capabilities in $Y_n$; see
Definition~\ref{def:category_partition}(i) applied internally to
$Y_n$), define the \emph{event-local} decomposition coefficients
$\alpha_{n,c} \geq 0$ for each capability $c \in Y_n$, satisfying:
\[
  \sum_{c \in Y_n} \alpha_{n,c} = 1
  \qquad\text{(partition of unity within the bundle)}.
\]
The coefficients represent the share of bundle $Y_n$ attributable
to capability~$c$, determined at the time of event~$n$ and
\emph{fixed thereafter}.  Each event carries its own
decomposition, independent of other events.  The decomposition can
be obtained from prior hedonic
regression~\citep{rosen1974hedonic} on similar trades, domain
knowledge, or the equal-weight default
$\alpha_{n,c} = 1/|Y_n|$.

Within-bundle poset-independence is the hypothesis that makes
the partition of unity compatible with S5's share inequality.
Under within-bundle independence, no cooperative node activates
from within $Y_n$, so
$\volR(Y_n) = \sum_{c \in Y_n} \volR(c)$ by axiom~M4, and S5's
$\alpha_{n,c} \leq \volR(c)/\volR(Y_n)$ sums to at most one
across $c \in Y_n$ with equality achievable.  Bundles with
internal cooperative activation fall back to Part~A only (see
Remark~\ref{rem:within_bundle_indep}).
\end{definition}

\begin{remark}[Why within-bundle independence is a Part~B
hypothesis]
\label{rem:within_bundle_indep}
If $Y_n$ contains capabilities $a, b$ whose cooperative
composition $z$ has all participants in $Y_n$, then S4(a)
(Section~\ref{sec:assumptions}) treats $z$ as an explicit
bundle summand, so
$\volR(Y_n) = \volR(a) + \volR(b) + \volR(z) > \sum_{c \in Y_n}
\volR(c)$.  S5's share inequality $\alpha_{n,c} \leq
\volR(c)/\volR(Y_n)$ is then forced below $\volR(c)/(\volR(a) +
\volR(b) + \volR(z))$ for each atomic $c \in Y_n$, and
$\sum_{c \in Y_n} \alpha_{n,c} \leq [\volR(a) +
\volR(b)]/\volR(Y_n) < 1$, contradicting the partition of
unity.  Part~B's per-capability attribution therefore requires
within-bundle poset-independence.  Part~A places no such
constraint: its bundle-level bound $\volR(Y_n) \geq
\Delta\volL(X_n)$ is valid regardless of internal cooperative
structure.
\end{remark}

\begin{remark}[Event-local versus global regression]
Classical hedonic regression~\citep{rosen1974hedonic} fits a
\emph{global} model across all observed trades simultaneously,
producing the tightest decomposition but breaking the monotonicity
of the max-attribution bound
(Proposition~\ref{prop:monotone}) when new observations refit all
coefficients.  The event-local formulation fixes coefficients at
observation time, sacrificing fit quality for the monotonicity
guarantee.
\end{remark}

\subsection{Central result}

\begin{theorem}[Aggregate B-to-C Lower Bound]
\label{thm:aggregate_bound}
Fix a trade window $W = [t_0, t_0 + T]$
(Definition~\ref{def:trade_window}) and a sequence of
revealed-sacrifice events
$\{(i_n, X_n, Y_n, t_n)\}_{n=1}^N$ with $t_n \in W$, whose
bundles jointly cover a subset $U$ of unbenchmarked
capability-space.  All $\volR$ quantities below refer to the
\emph{window-active} form $\volR^{[W]}$
(Definition~\ref{def:volR}): the theorem bounds the realized
volume accumulated during $W$, not the instantaneous
$\volR_t$ at any single evaluation time.

\medskip\noindent
\textbf{Part~A (Main result --- bundle-level lower bound;
requires S0--S2 and S4(a)).}
For each event~$n$:
\[
  \volR^{[W]}(Y_n) \;\geq\; \Delta\volL(X_n).
\]
For $N$ events with \emph{pairwise poset-independent} bundles
(in the sense of
Definition~\ref{def:category_partition}(i): no subsumption
relation and no cooperative-composition link connects any
capability in $Y_n$ to any capability in $Y_{n'}$ for
$n \neq n'$; this implies set-disjointness $Y_n \cap Y_{n'} =
\emptyset$ as a special case, since a capability in both
bundles would be trivially poset-related to itself):
\[
  \volR^{U,[W]} \;\geq\; \sum_{n=1}^{N} \Delta\volL(X_n),
\]
where $\volR^{U,[W]}$ is the window-active realized
capability volume restricted to the unbenchmarked subset $U$
covered by the union of bundles.  Part~A is the paper's
substantive result: the $\volL$-to-$\volR$ lower bound per unit
of observed sacrifice magnitude.

\medskip\noindent
\textbf{Part~B (Conditional refinement --- per-capability
attribution; requires S0--S5, within-bundle
poset-independence, and positive sacrifice magnitude).}
Restrict attention to events $n$ with $\Delta\volL(X_n) > 0$
whose bundles $Y_n$ are within-bundle poset-independent
(Definition~\ref{def:bundle_decomp}).  Events with
$\Delta\volL(X_n) = 0$ carry no Part-B information and are
excluded from the attribution.  When bundles overlap and
per-capability attribution is desired, define for each
capability $c \in U$:
\[
  \volRlower(c) \;=\; \max_{n:\, c \in Y_n}
  \alpha_{n,c} \cdot \Delta\volL(X_n).
\]
Then, provided the capabilities $c \in U$ are pairwise
poset-independent:
\[
  \volR^{U,[W]} \;\geq\; \sum_{c \in U} \volRlower(c).
\]
Part~B is \emph{conditional on S5}---a direct share inequality
on the quantity being bounded (see
Remark~\ref{rem:s5_heavy})---and should be read as a refinement
when admissible, not as a result comparable in strength to
Part~A.
\end{theorem}

\begin{proof}
\emph{Part~A.}
Fix an event $n$ with $t_n \in W$.  By S1, agent~$i_n$ chose
$Y_n$ over $X_n$ while both were available.  Revealed preference
gives $\Uin{n}\bigl((H_n \setminus X_n) \cup Y_n\bigr) \geq
\Uin{n}(H_n)$ for pre-trade holdings $H_n$.  By
Lemma~\ref{lem:cancellation} under S4(a) and the genuine-
exchange conditions of Definition~\ref{def:sacrifice_event},
this reduces to $\Uin{n}(Y_n) \geq \Uin{n}(X_n)$.  By S0
(valuation-bounded-by-exercise, stated in window-active
form), $\Uin{n}(Y_n) \leq \volR^{[W]}(Y_n)$.  By S2
(benchmark grounding), $\Uin{n}(X_n) \geq \Delta\volL(X_n)$.
Chaining:
\[
  \volR^{[W]}(Y_n) \;\geq\; \Uin{n}(Y_n)
  \;\geq\; \Uin{n}(X_n)
  \;\geq\; \Delta\volL(X_n),
\]
which is the per-event bound.  The window-active form is
essential: the sacrifice channel observes acquisition of $Y_n$
at $t_n$, not exercise at any particular instant.  S0's
consequent is therefore stated as a bound on
$\volR^{[W]}(Y_n)$---the cumulative exercise witness for $Y_n$
across $W$---not on the instantaneous $\volR_t(Y_n)$ at any
single evaluation time $t$, which could be zero even when
$Y_n$ was exercised earlier or later in $W$.

Suppose the $N$ bundles $\{Y_n\}$ are pairwise poset-independent:
no subsumption relation and no cooperative-composition link
connects any capability in $Y_n$ to any capability in $Y_{n'}$
for $n \neq n'$.  Then the sub-posets carried by the bundles are
disconnected components of the sub-poset on
$U = \bigcup_n Y_n$.  Axiom~M4
of~\citep[Proposition~1]{lasser2026poset} makes the poset
measure additive over poset-disjoint components; $\volR^{[W]}$
inherits M4 (Section~\ref{sec:axiom_inheritance},
Theorem~\ref{thm:axiom_inheritance}), so
$\volR^{U,[W]} = \sum_n \volR^{[W]}(Y_n)$.  Summing the
per-event bounds:
\[
  \volR^{U,[W]} \;=\; \sum_{n=1}^N \volR^{[W]}(Y_n)
  \;\geq\; \sum_{n=1}^N \Delta\volL(X_n),
\]
establishing Part~A.

\emph{Part~B.}
Part~B is stated for events with $\Delta\volL(X_n) > 0$ and
bundles $Y_n$ that are within-bundle poset-independent
(Definition~\ref{def:bundle_decomp}); events with zero sacrifice
carry no information, and bundles with internal cooperative
activation fall back to Part~A
(Remark~\ref{rem:within_bundle_indep}).  Fix an admissible event
$n$ and its bundle $Y_n$.  By Part~A, $\volR(Y_n) \geq
\Delta\volL(X_n) > 0$, so $\volR(Y_n) > 0$ and division is
well-defined.  S3 (bundle coherence) additionally ensures the
hedonic decomposition of $Y_n$ is non-degenerate, so
$\alpha_{n,c}$ is well-specified.

Within-bundle poset-independence plus~M4
(Section~\ref{sec:axiom_inheritance},
Theorem~\ref{thm:axiom_inheritance}) gives
$\volR(Y_n) = \sum_{c \in Y_n} \volR(c)$: no cooperative node
activates from within $Y_n$, so the bundle's $\volR$ decomposes
atomically.  S5 then asserts that the event-local coefficients
are lower bounds on the true atomic shares:
\[
  \alpha_{n,c} \;\leq\; \volR(c) / \volR(Y_n),
  \qquad c \in Y_n,
\]
which is feasible with $\sum_{c \in Y_n} \alpha_{n,c} = 1$
precisely under within-bundle independence.  Rearranging and
applying the Part~A per-event bound:
\[
  \volR(c) \;\geq\; \alpha_{n,c} \cdot \volR(Y_n)
  \;\geq\; \alpha_{n,c} \cdot \Delta\volL(X_n),
\]
for every $c \in Y_n$.  If $c$ appears in multiple admissible
bundles, each quantity $\alpha_{n,c} \cdot \Delta\volL(X_n)$ is a
valid lower bound on $\volR(c)$; taking the maximum preserves
validity:
\[
  \volRlower(c) \;:=\;
  \max_{n:\, c \in Y_n} \alpha_{n,c} \cdot \Delta\volL(X_n)
  \;\leq\; \volR(c).
\]
(Using the max rather than a sum is what underwrites the
monotonicity of Proposition~\ref{prop:monotone}: additional
events can only weakly increase a max.)

By the Part~B hypothesis, the capabilities $c \in U$ are
pairwise poset-independent, so a second application of~M4
(to the poset of capabilities, not bundles) gives global
additivity:
\[
  \volR^{U,[W]} \;=\; \sum_{c \in U} \volR(c)
  \;\geq\; \sum_{c \in U} \volRlower(c).
\]
\end{proof}

\begin{remark}[Why poset-independence, not set-disjointness]
\label{rem:poset_indep}
Set-disjointness ($Y_1 \cap Y_2 = \emptyset$) ensures no
\emph{capability} appears in two bundles, but does not rule out
subsumption or cooperative links across bundles.  If a cooperative
capability $c_{\mathrm{coop}}$ depends on capabilities in both
$Y_1$ and $Y_2$, its $\volR$ contribution is not attributable to
either bundle alone.  Poset-independence rules this out: under
poset-independence, the sub-posets are disconnected components
and $\volL$ (hence $\volR$) decomposes additively by
axiom~M4~\citep[Proposition~1]{lasser2026poset}.
\end{remark}

\begin{definition}[Subset and Category-Level $\volRlower$]
\label{def:volRlower_subset}
Under the Part~B hypotheses (S0--S5) of
Theorem~\ref{thm:aggregate_bound}, the per-capability lower
bound $\volRlower(c)$ extends to any subset $A \subseteq P$
whose covered capabilities are pairwise poset-independent:
\[
  \volRlower(A) \;:=\; \sum_{c \in A \cap U} \volRlower(c),
  \qquad \text{when } A \cap U \text{ is pairwise
  poset-independent,}
\]
where $U$ is the observed subset covered by bundles.  In this
regime the sum is a valid lower bound on $\volR(A)$, via
Part~B applied to the sub-family $A \cap U$.  When $A \cap U$
has internal poset-dependencies, the same expression can be
formed but is an \emph{accounting score}, not a
theorem-certified lower bound on $\volR(A)$.  In particular:
\begin{itemize}
\item For a bundle category $C_j \in \mathcal{C}$
  (Definition~\ref{def:category_partition}), $C_j \cap U$ is
  pairwise poset-independent only when the public partition is
  refined below internal dependencies.  For coarser partitions,
  $\volRlower(C_j)$ is an accounting score for the category.
\item For any other capability subset $A$ used by downstream
  diagnostics (restriction set, residual class, diagnostic
  cell), certified lower-bound status requires the same
  internal poset-independence of $A \cap U$.
\end{itemize}
Under Part~A alone, $\volRlower$ is defined only at the
per-bundle level; the global aggregate is $\volR^{U,[W]} \geq \sum_n
\Delta\volL(X_n)$, and arbitrary subset-level $\volRlower(A)$
is not a Part~A object.  The companion paper's gap decomposition
and wireheading diagnostic~\citep{lasser2026paper8b} are stated
under Part~B semantics and report subset/cell sums as
theorem-certified lower bounds only when the corresponding
internal poset-independence condition is met; otherwise the
sums are reported as diagnostic accounting scores.
\end{definition}

\subsection{B-to-C ratio and monotonicity}

\begin{definition}[B-to-C Ratio Under Revealed Sacrifice]
\label{def:btoc_ratio}
The B-to-C ratio under revealed-sacrifice observation is:
\[
  \betaL(G, T) \;=\; \volRlower(G, T) \;/\; \volL(G),
\]
where $G$ is a capability subset on which $\volRlower(G, T)$ is
defined.  The two reference choices are:
(a)~$G = U$, where $U$ is the observed subset covered by
admissible bundles (so $\volRlower(U, T)$ is the per-capability
or sub-batch-supremum aggregate of
Theorem~\ref{thm:aggregate_bound});
(b)~$G = P$, the full capability-space denominator, with
numerator $\volRlower(U, T)$ extended by zero outside $U$ to
yield a global conservative ratio.  In both cases
$\betaL \in [0, 1]$: the numerator is $\geq 0$ by construction
and is bounded above by $\volL(G)$ because $\volRlower(G, T)
\leq \volR(G) \leq \volL(G)$ (Part-A/Part-B inequalities chained
with $\volR \leq \volL$).  For arbitrary $G \subseteq P$ outside
case~(a) or~(b), $\volRlower(G, T)$ requires the subset-extension
hypotheses of Definition~\ref{def:volRlower_subset}; without
those, $\betaL(G, T)$ is an accounting score, not a
theorem-certified ratio.  $T$ indexes the trade window.
\end{definition}

\begin{proposition}[Part~B Monotone Accumulation Under
Max-Attribution]
\label{prop:monotone}
Under Part~B hypotheses (S0--S5 and within-bundle
poset-independence), define the per-capability lower bound as
the maximum attribution across all observed events:
\[
  \volRlower(c) \;=\; \max_{n:\, c \in Y_n}
  \alpha_{n,c} \cdot \Delta\volL(X_n),
\]
and the aggregate as $\volRlower = \sum_{c \in U} \volRlower(c)$.
Then the aggregate is non-decreasing in the number of observed
admissible events: for event sets $E \subseteq E'$ both satisfying
the Part~B admissibility conditions,
\[
  \volRlower(E') \;\geq\; \volRlower(E).
\]
\end{proposition}

\begin{proof}
Fix $c \in U$.  Let $S_N(c) = \{n \leq N : c \in Y_n\}$ be the
set of events in the first $N$ observations that cover $c$.  For
$N' > N$, $S_{N'}(c) \supseteq S_N(c)$, and
\[
  \volRlower^{(N')}(c)
  = \max_{n \in S_{N'}(c)} \alpha_{n,c} \Delta\volL(X_n)
  \;\geq\;
  \max_{n \in S_N(c)} \alpha_{n,c} \Delta\volL(X_n)
  = \volRlower^{(N)}(c),
\]
because a maximum over a superset is weakly larger.  If
$c \notin Y_n$ for any $n$ then both sides are zero by convention
and the inequality is trivial.  Summing the per-capability
inequalities over $c \in U$:
$\volRlower^{(N')} \geq \volRlower^{(N)}$.
\end{proof}

\begin{proposition}[Part~A Monotone Accumulation via
Sub-Batch Supremum]
\label{prop:monotone_part_a}
Under Part~A hypotheses (S0--S2, S4(a), and the
genuine-exchange conditions of
Definition~\ref{def:sacrifice_event}), define the Part~A lower
bound as the supremum over admissible pairwise-poset-independent
sub-batches of the observed event set:
\[
  \volRlower^{\mathrm{A}}(E) \;:=\;
  \sup\Bigl\{\, \sum_{n \in B} \Delta\volL(X_n)
  \;:\; B \subseteq E \text{ and } B \text{ satisfies the
  Part~A admissibility and pairwise poset-independence
  hypotheses of Theorem~\ref{thm:aggregate_bound}} \,\Bigr\},
\]
with the supremum attaining $0$ on the empty sub-batch.  Then
$\volRlower^{\mathrm{A}}$ is non-decreasing in $E$: for event
sets $E \subseteq E'$,
\[
  \volRlower^{\mathrm{A}}(E') \;\geq\;
  \volRlower^{\mathrm{A}}(E).
\]
\end{proposition}

\begin{proof}
Any admissible sub-batch $B \subseteq E$ is also an admissible
sub-batch $B \subseteq E'$ (admissibility is a property of $B$
itself, not of its ambient event set).  Hence the supremum over
admissible sub-batches of $E$ is a supremum over a subset of the
admissible sub-batches of $E'$, and the supremum over a superset
is weakly larger.
\end{proof}

\begin{remark}[Interpretation of Part~A monotonicity]
\label{rem:part_a_monotone_interp}
$\volRlower^{\mathrm{A}}(E)$ is the tightest Part~A lower bound
the framework can extract from the event set $E$ by choosing
which sub-batch to apply Theorem~\ref{thm:aggregate_bound} to.
The supremum is over single sub-batches; computing it does
\emph{not} entail summing across multiple disjoint sub-batches
of $E$, because doing so could double-count when an underlying
capability appears in multiple sub-batches.  Operationally, the
framework selects the single sub-batch maximising the Part~A
sum; if multiple maximal admissible sub-batches partition the
covered space into poset-disjoint components, their sum is also
admissible (and equal to the supremum on that one combined
sub-batch), but the principled definition is the supremum over
sub-batches, not the sum across them.  $\volRlower^{\mathrm{A}}$
is what the paper refers to when it says ``the constructed
Part~A lower bound is non-decreasing in observed admissible
events.''
\end{remark}
